\section*{Chapter \ref{ch:ml:interpretability-and-model-governance} --- Interpretability and Model Governance}

\begin{solution}{exo:ml:interpretability-and-model-governance:1}
$\mathrm{PD}(v) = v + \E[x_2^2] = v + 1/3$: the true effect plus a constant. If $x_2 = x_1$ and the model is $f$ itself,
\omterm{def:ml:interpretability-and-model-governance:pd}{partial dependence} still averages $x_2^2$ over the marginal of $x_2$ and returns $v + 1/3$, while along the data the
prediction is $v + v^2$: \omterm{def:ml:interpretability-and-model-governance:pd}{partial dependence} describes the formula, not what the model does on the data, and attributes none
of the curvature to the only variable that moves.
\end{solution}

\begin{solution}{exo:ml:interpretability-and-model-governance:2}
$0.1\times250 = 25$ expected. The probability of 27 or more is 0.37: not alarming.
\end{solution}

\begin{solution}{exo:ml:interpretability-and-model-governance:3}
The tree's rules describe the part of the model's behaviour it reproduces, less than half of the variance here; nothing may be
concluded about the rest, which here is the whole interaction. A surrogate's rules are an explanation only where its fidelity
is high, and should be checked on the slices that matter.
\end{solution}

\begin{solution}{exo:ml:interpretability-and-model-governance:4}
Its effect is $2x_3x_4$ with $x_4$ symmetric around zero: for every observation where raising $x_3$ raises the prediction there
is one where it lowers it, and the average slope is zero. The ICE curves' slopes follow $x_4$.
\end{solution}

\begin{solution}{exo:ml:interpretability-and-model-governance:5}
The five weak features have equal true effects; their estimated importances differ by sampling noise, which reorders them in
every resample. Only differences of importance larger than that noise are stable.
\end{solution}

\begin{solution}{exo:ml:interpretability-and-model-governance:6}
A share of SHAP attribution from one fit, with correlated features and no stability check, does not say the model is a
revisions model: correlated features share credit arbitrarily, the share may change at the next retrain, and attribution is
not causation. Show ALE and ICE for the feature, the attribution across retrains, and the performance of the model without it.
\end{solution}

\begin{solution}{exo:ml:interpretability-and-model-governance:7}
It appears to rise to perfection: the top five are features 0 to 4 in every retrain and the Kendall correlation is 1. The
explanation is that nothing was learned: with 400 observations, \omterm{def:ml:trees-and-boosting:bagging}{bagging} at 70\% and the firm's minimum of 200 observations per
leaf, no tree can split, every \omterm{def:ml:feature-engineering-selection-and-importance:shapley}{SHAP value} is zero, and the ranking falls back to the column order. A stability score must be read
with the importances it ranks.
\end{solution}

\begin{solution}{exo:ml:interpretability-and-model-governance:8}
ALE of $j$ accumulates $\E[f(z_{k}, x_{-j}) - f(z_{k-1}, x_{-j})\mid x_j\in(z_{k-1}, z_k]]$; for an additive $f$ the other
components cancel in each difference, leaving $f_j(z_k) - f_j(z_{k-1})$, and the sum telescopes to $f_j$ plus a constant.
\omterm{def:ml:interpretability-and-model-governance:pd}{Partial dependence} is $f_j(v) + \sum_{k\ne j}\E f_k(x_k)$, also $f_j$ plus a constant. Both are right for the true additive
function; \omterm{def:ml:interpretability-and-model-governance:pd}{partial dependence} goes wrong for a fitted model, which is only constrained where the data are and may not be
additive off the data's support, where \omterm{def:ml:interpretability-and-model-governance:pd}{partial dependence} evaluates it.
\end{solution}

\begin{solution}{pb:ml:interpretability-and-model-governance:1}
\textbf{Part I.}
\begin{enumerate}
\item At correlations 0, 0.92 and 0.99: PD 0.027, 0.134, 0.229; ALE 0.028, 0.044, 0.110.
\item The fitted model itself cannot separate nearly identical features; ALE explains the model, which mixes their effects.
\item Slopes spread with standard deviation 0.91 and correlation 0.98 with $x_4$: the interaction.
\item The additive part in $x_1$ and $x_2$ (fidelity 0.48); none of the interaction.
\end{enumerate}
\textbf{Part II.}
\begin{enumerate}[resume]
\item The prediction (0.82) comes mostly from the interaction; $x_1$ would have to fall from 0.46 to $-0.44$ to reverse it.
\item At $R^2$ 0.66: the same set in 22\% of retrains, Kendall 0.89; at $R^2$ 0.05: never, Kendall 0.78, with null features
sometimes in the top five.
\item Importances from several retrains, their spread, the features stable at the top, and a warning on correlated groups.
\item With the deployed model's \omterm{def:ml:feature-engineering-selection-and-importance:shapley}{SHAP values} for that stock on Monday and Tuesday, the inputs that changed, a counterfactual,
and the ALE of those inputs showing that the response is sensible.
\end{enumerate}
\textbf{Part III.}
\begin{enumerate}[resume]
\item Fifteen fields from name to approvals; monitoring and approvals are missing.
\item Out-of-sample performance ($R^2$ 0.860), the challenger (ridge 0.349), implementation parity (exact on 200 rows) and
outcomes (27 of 250, tail 0.37).
\item The model that runs must be the model that was validated; exports, compilers (chapter~26) and retrains break that silently.
\item That supervisors treat learned models as models, under the same principles, whatever the technology.
\end{enumerate}
\textbf{Part IV.}
\begin{enumerate}[resume]
\item \emph{Explain it to the risk committee.} \omterm{def:ml:interpretability-and-model-governance:pd}{Partial dependence}'s error against the true effect is 0.134 and 0.229 at
correlations 0.92 and 0.99, against ALE's 0.044 and 0.110; the top five SHAP features are the same set in 22\% of ten retrains
at $R^2$ 0.66 (Kendall 0.89) and in none at $R^2$ 0.05 (Kendall 0.78).
\item Not here, where the interaction carries half the signal and the challenger loses it; yes where a linear model is within
noise of the boosted one.
\item A high tier if it drives positions at size: material, complex and uncertain.
\item An independent team rebuilding the model from the card, a challenger, stress inputs outside the training range, and
explanations checked against economic sense.
\item Performance, stability, explanations and parity; the card's version and data fields.
\item At the level of the features that feed the embedding, or with probes that relate embedding directions to known features.
\item Secrets (credentials, client data) and claims without evidence.
\item A model of the model, with an error of its own.
\end{enumerate}
\end{solution}

\begin{solution}{iq:ml:interpretability-and-model-governance:1}
\omterm{def:ml:interpretability-and-model-governance:pd}{Partial dependence} averages predictions with the feature forced to each value for every observation, visiting combinations
absent from the data; ALE averages local differences within intervals, for observations actually there. With correlated
features \omterm{def:ml:interpretability-and-model-governance:pd}{partial dependence} can be badly wrong; ALE stays on the data.
\iqlookfor{the two definitions and the off-manifold problem.}
\end{solution}

\begin{solution}{iq:ml:interpretability-and-model-governance:2}
Often unstable below the top few features, especially with weak signals and correlated features. Measure by retraining on
resamples or successive periods and reporting top-set agreement and rank correlations, alongside the importances themselves.
\iqlookfor{instability of the tail of the ranking and a measurement protocol.}
\end{solution}

\begin{solution}{iq:ml:interpretability-and-model-governance:3}
Purpose and use, target and horizon, data and period, features, training procedure, validation results with a challenger,
performance by regime, explanations, limitations, monitoring, owners and approvals, version.
\iqlookfor{purpose, data, evaluation, limitations and ownership.}
\end{solution}

\begin{solution}{iq:ml:interpretability-and-model-governance:4}
The forecast's change and the inputs that drove it (SHAP on the deployed model), whether the response is sensible (ALE/ICE),
what would have reversed it, and how the position rule turned the forecast into the trade.
\iqlookfor{local attribution, sanity of the response, and the link to the position.}
\end{solution}

\begin{solution}{iq:ml:interpretability-and-model-governance:5}
Out-of-time performance against a simple challenger, leakage checks, stability across retrains, explanations checked for sign
and size, implementation parity of the exported model, outcomes analysis, a \omterm{def:ml:interpretability-and-model-governance:card}{model card}, and monitoring before go-live.
\iqlookfor{challenger, leakage, stability, parity and monitoring.}
\end{solution}

\begin{solution}{iq:ml:interpretability-and-model-governance:6}
When the interpretable model is within noise of the black box, when decisions must be justified case by case, and when the
cost of an unexplained failure is high; the weak signals of trading often make the first condition true.
\iqlookfor{performance parity and the cost of opacity.}
\end{solution}
